\section*{Capítulo \ref{ch:g9:landforms} --- Las formas del relieve}

\begin{solution}{exo:g9:landforms:1}
(a) El viento. (b) El agua que disuelve la \omterm{def:g3:stones-pebbles-sand:rock}{roca} (y que deposita calcita en
la \omterm{def:g9:landforms:karst}{cueva}). (c) El \omterm{def:g2:water-on-earth:ocean}{mar} (\omterm{def:g9:landforms:beach}{deriva litoral}). (d) El hielo. (e) El \omterm{def:g2:water-on-earth:ocean}{mar}. (f) Un
\omterm{def:g1:landscapes-around-us:water}{río}.
\end{solution}

\begin{solution}{exo:g9:landforms:2}
(a) El \omterm{def:g9:landforms:uplift}{levantamiento}. (b) La \omterm{def:g9:landforms:dune}{cara de avalancha}. (c) Una \omterm{def:g9:landforms:karst}{dolina}. (d) La
\omterm{def:g9:landforms:beach}{deriva litoral}. (e) Un \omterm{def:g9:landforms:karst}{karst}.
\end{solution}

\begin{solution}{exo:g9:landforms:3}
(a) Falso: los cuernos apuntan en el sentido del viento; el frente
redondeado mira al viento. (b) Verdadero. (c) Falso: el viento la desplaza
grano a grano, ladera arriba por la pendiente suave y abajo por la \omterm{def:g9:landforms:dune}{cara de
avalancha}. (d) Verdadero.
\end{solution}

\begin{solution}{exo:g9:landforms:4}
(b), (d), (e), (a), (c).
\end{solution}

\begin{solution}{exo:g9:landforms:5}
(a) Del oeste: los cuernos apuntan en el sentido del viento. (b) En el
lado este, entre los cuernos. (c) $60 \div 6 = 10$ años. (d) Una \omterm{def:g9:landforms:dune}{duna} más
grande contiene más \omterm{def:g3:stones-pebbles-sand:sand}{arena}, y toda ella tiene que pasar por encima de la
cresta para que la \omterm{def:g9:landforms:dune}{duna} avance.
\end{solution}

\begin{solution}{exo:g9:landforms:6}
(a) $7 \times 1 = \qty{7}{m}$. (b) $2 \div 12 = \frac{1}{6}$ de año.
(c) Siete años de retroceso en $\frac{1}{6}$ de año: $7 \times 6 = 42$
veces más deprisa. (d) Las olas excavan primero una muesca; la \omterm{def:g3:stones-pebbles-sand:rock}{roca} de
encima queda un tiempo en voladizo, sin moverse; después el voladizo se
derrumba de golpe, y el borde retrocede de un salto.
\end{solution}

\begin{solution}{exo:g9:landforms:7}
(a) Hacia el este. (b) Tres trazos \omterm{def:g9:landforms:beach}{playa} arriba inclinados hacia el
noreste, cada uno seguido de un trazo de vuelta en línea recta hacia el
sur: un zigzag que avanza hacia el este. (c) Una \omterm{def:g9:landforms:beach}{flecha litoral}, que crece
hacia el este a través de la entrada de la bahía. (d) La flecha protege de
las olas el agua que queda detrás; en el agua tranquila, el \omterm{def:g3:stones-pebbles-sand:sand}{barro} fino
puede depositarse.
\end{solution}

\begin{solution}{exo:g9:landforms:8}
(a) Absorbe dióxido de carbono del aire y del suelo. (b) El ácido débil
disuelve despacio la caliza, pero apenas ataca el granito. (c) El agua
solo puede entrar en la \omterm{def:g3:stones-pebbles-sand:rock}{roca} por sus grietas y a lo largo de sus capas;
ensancha primero esas aberturas. (d) La \omterm{def:g3:stones-pebbles-sand:rock}{roca} que hay bajo la superficie se
disuelve y el terreno se hunde lentamente; o se derrumba el techo de una
\omterm{def:g9:landforms:karst}{cueva}.
\end{solution}

\begin{solution}{exo:g9:landforms:9}
(a) Cuando la gota se encuentra con el aire de la \omterm{def:g9:landforms:karst}{cueva}, desprende parte de
su dióxido de carbono y ya no puede retener toda la caliza disuelta, que
sale en forma de calcita. (b) Una estalagmita crece a partir de gotas que
salpican el suelo y se extienden; no hay ninguna gota colgando de su punta
que construya un tubo. (c) $\qty{30}{cm} = \qty{300}{mm}$; $300 \div 0.1 = 3000 = 3 \times 10^3$
años.
\end{solution}

\begin{solution}{exo:g9:landforms:10}
(a) $4600 \div 10^7 = 4.6 \times 10^{-4}$~m al año, o \qty{0.46}{mm} al
año. (b) $0.46 - 0.2 = \qty{0.26}{mm}$ al año. (c) $\qty{1}{m} =
\qty{1000}{mm}$; $1000 \div 0.46 \approx 2200$, unas dos mil veces más
rápido. (d) Se mantiene durante millones de años: $\qty{0.26}{mm} \times
10^7 = 2.6 \times 10^6$~mm, unos \qty{2.6}{km}.
\end{solution}

\begin{solution}{exo:g9:landforms:11}
(a) Forma: \omterm{def:g9:landforms:dune}{dunas} en media luna con su \omterm{def:g9:landforms:dune}{cara de avalancha} en el lado
sureste. Agente y proceso: el viento, que sopla del noroeste;
\omterm{def:g5:water-shapes-the-land:deposition}{sedimentación}. Material: \omterm{def:g3:stones-pebbles-sand:sand}{arena} suelta. (b) Forma: un \omterm{def:g9:landforms:cliff}{acantilado} sobre una
superficie plana descubierta con la marea baja. Agente y proceso: el \omterm{def:g2:water-on-earth:ocean}{mar},
\omterm{def:g5:water-shapes-the-land:erosion}{erosión}; la superficie es una plataforma de abrasión que ha dejado el
\omterm{def:g9:landforms:cliff}{acantilado} al retroceder. Material: \omterm{def:g3:stones-pebbles-sand:rock}{roca} blanca, probablemente creta, una
\omterm{def:g3:stones-pebbles-sand:rock}{roca} blanda. (c) Forma: un arroyo que desaparece por un agujero, un
\omterm{def:g6:water-underground:spring}{manantial} lejos de allí. Agente y proceso: el agua que disuelve la \omterm{def:g3:stones-pebbles-sand:rock}{roca} y
circula bajo tierra. Material: caliza; la meseta es un \omterm{def:g9:landforms:karst}{karst}.
\end{solution}

\begin{solution}{exo:g9:landforms:12}
(a) La \omterm{def:g9:landforms:dune}{duna} pequeña: $45 \div 6 = 7.5$ años; la más grande: $45 \div 3 =
15$ años. (b) $45 \div 100 = 0.45$ de año, unos cinco o seis meses.
(c) Por ejemplo, plantar hierbas en las \omterm{def:g9:landforms:dune}{dunas}: sus raíces sujetan la
\omterm{def:g3:stones-pebbles-sand:sand}{arena}, así que el viento ya no puede llevársela.
\end{solution}

\begin{solution}{pb:g9:landforms:1}
\textbf{1.} $46.0 - 30.4 = \qty{15.6}{m}$.

\textbf{2.} $15.6 \div 20 = \qty{0.78}{m}$ al año.

\textbf{3.} Entre 2012 y 2014: $41.0 - 36.2 = \qty{4.8}{m}$ en dos
años, es decir, $4.8 \div 2 = \qty{2.4}{m}$ al año.

\textbf{4.} $2.4 \div 0.78 \approx 3$ veces más deprisa. Probablemente un
invierno tormentoso, durante el cual se derrumbaron grandes voladizos.

\textbf{5.} \qty{78}{cm} al año; $0.78 \times 100 = \qty{78}{m}$ por
siglo.

\textbf{6.} Sí: \qty{0.78}{m} al año es algo menos de un metro al año.

\textbf{7.} $17 \div 0.78 \approx 22$ años.

\textbf{8.} El \omterm{def:g9:landforms:cliff}{acantilado} retrocede a saltos: un intervalo de dos años
puede contener un gran desprendimiento, o ninguno. A lo largo de 20 años,
los saltos y los años tranquilos se compensan.

\textbf{9.} Al borde le quedan $30.4 - 10 = \qty{20.4}{m}$: $20.4 \div 0.78
\approx 26$ años, así que en 2050.

\textbf{10.} $2 \times 0.78 = \qty{1.56}{m}$ al año; $30.4 \div 1.56
\approx 19.5$ años después de 2024, así que en 2043.

\textbf{11.} El borde de un \omterm{def:g9:landforms:cliff}{acantilado} puede ceder en cualquier momento;
situarse sobre él, o cerca de él, es peligroso.

\textbf{12.} $30.4 \div 0.78 \approx 39$ años después de 2024: en
\textbf{2063}.
\end{solution}
